\section{COMANDOS TERMINAL \normalfont{(\texttt{Wnd+R})}}
\newcommand{\kbraw}[4]{#1 & \texttt{#2} & #3 & \texttt{#4} \\}


\noindent
\begin{tabularx}{\linewidth}{@{} X l | X l @{}}
\multicolumn{2}{c}{\texttt{SISTEMA Y EJECUCIÓN}} & \multicolumn{2}{c}{\texttt{ARCHIVOS Y DIRECTORIOS}} \\
\midrule
\kbraw{Listar archivos}{ls}{Crear directorio}{mkdir [nombre]}
\kbraw{Ruta actual}{pwd}{Entrar a carpeta}{cd [ruta]}
\kbraw{Procesos activos}{Get-Process}{Subir un nivel}{cd ..}
\kbraw{Limpiar pantalla}{cls}{Borrar (recursivo)}{rm -rf [nombre]}
\kbraw{Editar archivo}{code [archivo]}{Copiar (recursivo)}{cp -r [orig] [dest]}
\kbraw{Manual/Ayuda}{help [comando]}{Mover o Renombrar}{mv [viejo] [nuevo]}
\kbraw{Ejecutar script}{python script.py}{Ver contenido}{cat [archivo]}
\midrule
\end{tabularx}
\noindent
\begin{tabularx}{\linewidth}{@{} X l | X l @{}}
\multicolumn{2}{c}{\texttt{ENTORNO PYTHON (PIP)}} & \multicolumn{2}{c}{\texttt{UTILIDADES DE RED/DATOS}} \\
\midrule
\kbraw{Instalar paquete}{pip install [paquete]}{Mostrar estructura}{Tree}
\kbraw{Lista paquetes}{pip list}{Descargar archivo}{wget [url] -OutFile [n]}
\kbraw{Instalar requisitos}{pip install -r req.txt}{Conexión Remota}{ssh usuario@ip}
\kbraw{Actualizar pip}{pip install -U pip}{Copiar a remoto}{scp [f] user@ip:[path]}
\bottomrule
\end{tabularx}



\section{SHORTCUTS VS Code \normalfont{(\texttt{Ctrl+K Ctrl+S})}} 
% Comando para facilitar la edición: \kbrow{Acción Izq}{Atajo Izq}{Acción Der}{Atajo Der}
\newcommand{\kbrow}[4]{#1 & \texttt{#2} & #3 & \texttt{#4} \\}

\begin{tabularx}{\linewidth}{@{} X l | X l @{}}
\multicolumn{2}{c}{\texttt{INTERFAZ}} & \multicolumn{2}{c}{\texttt{EDICIÓN}} \\
\midrule
\kbrow{Ajustes}{Ctrl+,}{Comentar}{Ctrl+Ç}
\kbrow{Show/Hide Sidebar}{Ctrl+B}{Bloque comentario}{Shift+Alt+A}
\kbrow{Show/Hide Second Sidebar}{Ctrl+Alt+B}{Select all}{Ctrl+A}
\kbrow{Quick Open}{Ctrl+P}{Select ocurrencia}{Ctrl+D}       
\kbrow{Paleta comandos}{Ctrl+Shift+P}{Select línea}{Ctrl+L}   
\kbrow{Abrir Explorador}{Ctrl++Shift+E}{Copiar línea}{Shift+Alt+$\uparrow\!\!/\!\!\downarrow$}
\kbrow{Búsqueda Global}{Ctrl+Shift+F}{Mover línea}{Alt+$\uparrow\!\!/\!\!\downarrow$}
\kbrow{Terminal}{Ctrl+Ñ}{Insertar línea abajo}{Ctrl+Shift+Enter}
\kbrow{Dividir/Unir editor actual}{Ctrl+Alt+$\rightleftarrows$}{Borrar línea}{Ctrl+Shift+K}
\kbrow{Dividir Editor}{Ctrl+\textbackslash}{Multicursor}{Ctrl+Alt+$\uparrow\!\!/\!\!\downarrow$}
\midrule
\addlinespace[2ex]
\multicolumn{2}{c}{\texttt{NAVEGACIÓN}} & \multicolumn{2}{c}{\texttt{ARCHIVOS}} \\
\midrule
\kbrow{Navegar editor}{Ctrl+Re/AvPag}{Nuevo archivo}{Ctrl+N}
\kbrow{Vista previa definición}{Alt+F12}{Guardar}{Ctrl+S}
\kbrow{Ir a definición}{Ctrl+Click}{Guardar todo}{Ctrl+K S}
\kbrow{Ir a símbolo}{Ctrl+Shift+O}{Cerrar pestaña}{Ctrl+W}
\kbrow{Buscar en archivo}{Ctrl+F}{Reabrir pestaña}{Ctrl+Shift+T}
\kbrow{Reemplazar}{Ctrl+H}{Cerrar carpeta}{Ctrl+K F}
\kbrow{Ir a línea}{Ctrl+G}{
Correr Python (propio)}{Ctrl+R}
\kbrow{Show Problems}{Ctrl+Shift+M}{Cambiar lenguaje}{Ctrl+K M}
\bottomrule
\end{tabularx}





\section{INSTALACIÓN WINDOWS DE PYTHON Y PIP}
\textbf{1.}: Descarga desde \href{https://www.python.org/}{\texttt{python.org}} $\rightarrow$ \text{Check: Add Python to PATH}.\\
\textbf{2. }: \texttt{python -m pip --version} \textit{(Si falta: \texttt{python -m ensurepip --default-pip})}.
\hyperlink{pip_venv}{Continuación en venv.}\\



\section{CONFIGURACIÓN BASE VS CODE}
\textbf{Extensiones}: 
\texttt{ms-python.python},
\texttt{ms-toolsai.jupyter}

\textbf{Ajustes}: \texttt{Auto-Save:\,After Delay, Format on save:\,\checkmark, Python Language Server:\,Pylance, Line Numbers:\,\checkmark}


\section{GESTIÓN DE ENTORNOS (venv) -- \normalfont{\scalebox{0.8}{\text{desde VS CODE}}}}
Entorno virtual para cada proyecto: aisla dependencias y evita problemas de compatibilidad.\\
Créalo al iniciar el proyecto, abre cuando estés trabajando y cierra al finalizar de trabajar.\\
\textbf{Crear}: \texttt{Abre dir proyecto → Ctrl+Ñ → python -m venv .venv → Ctrl+Shift+P → Python Interpreter → .venv}\\
\hypertarget{pip_venv}{
\text{Actualizar Pip Local}}: \texttt{python -m pip install --upgrade pip} \textit{(Ejecutar siempre tras activar el venv)}\\
\text{Verificar Ruta}: \texttt{where pip} debe apuntar a la carpeta \texttt{.venv}\\
\text{Tras crear el entorno virtual, configura lo siguiente}:\\
\text{Formato}: \texttt{pip install black → Ctrl+, → Python Formattixng Provider → black}\\
\text{Instala básicos}: \texttt{pip install matplotlib numpy scipy pandas scikit-learn jupyter}\\
\text{Persistencia}: \texttt{pip freeze > requirements.txt} \textit{(Documenta versiones)}\\
\textbf{Activar}: \texttt{Ctrl+Ñ → .venv\textbackslash Scripts\textbackslash activate} \;\;\; \textbf{Desactivar}: \texttt{Ctrl+Ñ → deactivate}\\ 
\text{Replicar entorno}: \texttt{pip install -r requirements.txt} \textit{(En un nuevo PC/clon)}


\section{CONTROL DE VERSIONES (GIT)}
\textbf{Configuración Inicial}: Establece tu identidad dentro del ecosistema \texttt{Git}\\
\texttt{git config --global user.name "Tu Nombre"} \\
\texttt{git config --global user.email "tu@email.com"}\\
\texttt{git config --list}: Lista tu configuración de nombre y correo.
\\
\textbf{Flujo de Trabajo Local}\\
\texttt{git init}: Inicializar repositorio.\\
\texttt{git status}: Listar archivos modificados y estado del índice.\\
\texttt{git diff}: Mostrar la diferencia ($\Delta$) entre el directorio de trabajo y el índice.\\
\texttt{git add .}: Indexar todos los cambios.\\
\texttt{git commit -m "mensaje"}: Crear un snapshot (punto de control) en el historial.\\
\texttt{git status}: Muestra el estado del directorio de trabajo y del área de preparación. \\
\texttt{git log --oneline --graph}: Visualización del grafo de evolución del proyecto.
\\\textit{Nota:} Si aparecen files de directorios padres, has inicializado (\texttt{git init}) en el nodo incorrecto. Antes de ejecutar \texttt{Remove-Item -Recurse -Force .git}, verifica ruta con \texttt{git rev-parse --show-toplevel}.\\
\textbf{Ramificaciones (Branching)}\\
\texttt{git branch <nombre>}: Crear una rama para experimentación aislada.\\
\texttt{git checkout <name>}: Cambiar a una rama.
\\\texttt{git checkout -b <nombre>}: Crear y cambiar a la nueva rama simultáneamente.\\
\texttt{git branch -a}: Lista todas las ramas del repositorio, tanto las locales como las remotas.\\
\texttt{git merge <nombre>}: Fusionar los cambios de la rama especificada en la rama activa.\\
\textbf{Sincronización Remota}\\
\texttt{git remote add origin <url servidor>}: Vincular repositorio local con servidor (GitHub).\\
\texttt{git push -u origin main}: Cargar commits locales al servidor remoto.\\
\texttt{git pull}: Descargar y fusionar cambios del remoto (equivale a \texttt{fetch} + \texttt{merge}).\\ 
\textbf{Gestión de Desastres y Retroceso (Undo)}\\
\texttt{git restore <archivo>}: Descartar cambios locales (volver al último commit).\\
\texttt{git reset --soft HEAD$\sim$1}: Deshacer último commit manteniendo cambios en \textit{staging}.\\
\texttt{git reset --hard HEAD$\sim$1}: \textbf{Peligro.} Borra commit y cambios no guardados.\\
\texttt{git restore --staged <archivo>}: Sacar del índice sin modificar contenido.\\
\textbf{Colaboración y Resolución de Conflictos}\\
\texttt{git fetch}: Descargar metadatos del remoto sin modificar archivos locales.\\
\texttt{git merge <rama>}: Fusionar. Si hay \textbf{Conflictos}, editar manualmente, \texttt{git add} y \texttt{git commit}.\\
\texttt{git pull}: Atajo de \texttt{fetch} + \texttt{merge}. Riesgo de conflictos directos.\\
\textbf{Organización y Depuración}\\
\texttt{git stash}: Almacenamiento temporal de cambios no commiteados (limpieza rápida).\\
\texttt{git show <hash>}: Inspeccionar cambios detallados de un commit específico.\\
\texttt{git blame <archivo>}: Listar autor y fecha de modificación por cada línea del archivo.


{
\setlength{\columnseprule}{0pt}
\begin{multicols}{2}
    \section*{ESTRUCTURA DE PROYECTO}
    \begin{Verbatim}
project/
├── .venv/         # Entorno local
├── .gitignore     # Filtro de Git
├── data/
│   ├── raw/       # Datos inmutables
│   └── processed/
├── notebooks/     # (.ipynb) 
├── src/           # (.py)
│   ├── __init__.py
│   └── tools.py
└── README.md
\end{Verbatim}
\columnbreak
\vspace*{-7pt}
    \begin{Verbatim}
    FICHERO .gitignore
__pycache__/
*.pyc
*$py.class
.venv/
build/
dist/
.ipynb_checkpoints
.vscode/
*.code-workspace
\end{Verbatim}
\end{multicols}
}

-- \texttt{.ipybn, .gitignore}: Los creas desde el explorador, no con \texttt{Ctrl+N}. Kernel debe apuntar a .venv\\



\section{ESTRUCTURA DE CELDAS (\texttt{ipython})}
\textbf{Marcador de celda}: \texttt{\# \%\%} (convierte \texttt{.py} en entorno modular).\\
\textbf{Ejecución rápida}: \texttt{Shift + Enter} para enviar bloque al núcleo interactivo.\\
\textbf{Plot Viewer}: Icono de imagen en gráficas renderizadas. Zoom y exportación directa.\\
-- \texttt{.py}: Outputs temporales, comentarios de código estándar, ejecución directa como script.\\
-- \texttt{.ipynb}: Outputs persistentes, Markdown y \LaTeX, requiere conversión (\texttt{nbconvert}).

